\documentclass[11pt]{article}
\usepackage[margin=1in]{geometry}
\usepackage{amsmath,amssymb,amsthm,mathtools}
\usepackage{graphicx}
\usepackage{microtype}
\usepackage{xcolor}
\usepackage[colorlinks=true,linkcolor=blue!60!black,citecolor=blue!60!black,urlcolor=blue!60!black]{hyperref}

\newtheorem{theorem}{Theorem}
\newtheorem{lemma}{Lemma}
\newtheorem{remark}{Remark}
\newcommand{\ZF}{\mathsf{ZF}}
\newcommand{\DC}{\mathsf{DC}}
\newcommand{\BP}{\mathsf{BP}}
\newcommand{\R}{\mathbb R}
\newcommand{\Sph}{S^{n-1}}
\newcommand{\C}{C(\Sph)}
\newcommand{\Vt}{\widetilde V}

\title{A choice-free model for the dual mixed volume characterization}
\author{Solution packet for arXiv:1408.6796, Question 3.5}
\date{11 August 2026}

\begin{document}
\maketitle

\begin{abstract}
We answer Question 3.5 of Jim\'enez--Villanueva affirmatively in relative
consistency form.  In every model of $\ZF+\DC$ in which all sets of reals have
the Baire property, separate radial additivity already forces the continuity
missing from the source theorem.  Consequently the asserted characterization
by dual mixed volume holds in every dimension.  Shelah's 1984 consistency
theorem supplies such a model from $\operatorname{Con}(\ZF)$, so no
inaccessible-cardinal assumption is required.
\end{abstract}

\section{The question and the answer}

Write $\mathcal S_0^n$ for the star bodies in $\R^n$, with radial addition,
and recall
\[
  \Vt(L_1,\ldots,L_n)
   =\frac1n\int_{\Sph}\rho_{L_1}(u)\cdots\rho_{L_n}(u)\,d\sigma(u).
\]
The question reproduced below asks whether it is compatible with $\ZF$ that
separate additivity, rotation invariance, and vanishing whenever two arguments
meet only at the origin characterize the scalar multiples of $\Vt$.

\begin{center}
\includegraphics[width=.94\textwidth]{figures/open_question.png}
\end{center}

For an integer $n\ge2$, let $D_n$ denote the assertion that every map
\[
 F:(\mathcal S_0^n)^n\longrightarrow\R
\]
having those three properties satisfies $F=c\Vt$ for some $c\in\R$.

\begin{theorem}[Full relative-consistency answer]\label{thm:main}
If $\ZF$ is consistent, then so is
\[
 \ZF+\text{``$D_n$ holds for every integer $n\ge2$''}.
\]
More concretely, every model of
\[
 \ZF+\DC+\text{``every set of reals has the Baire property''}
 \tag{$\ast$}
\]
satisfies $D_n$ simultaneously for all $n\ge2$.
\end{theorem}

\subsection*{Idea of the proof}
Radial functions identify star bodies with the positive cone of the separable
Banach space $C(\Sph)$.  Extend the given functional additively to the whole
space.  In an all-sets-Baire-property model, each one-variable section is a
Baire-measurable homomorphism between Polish groups, hence continuous by
Pettis' theorem.  This supplies positive homogeneity and exactly the separate
boundedness missing from the question.  The bounded characterization already
proved in the source paper then gives integration against a measure on the
sphere, and rotation invariance makes that measure spherical Lebesgue measure
up to scale.

The proof separates the set-theoretic input from the functional-analytic
argument.  The key point is that the sections of $F$ live not merely on
$\R$, but on the infinite-dimensional Polish group $\C$.

\section{Automatic continuity}

We use two standard consequences of $(\ast)$.

\begin{lemma}[Baire transfer and Pettis continuity]\label{lem:pettis}
Assume $(\ast)$.  Then every subset of any Polish space has the Baire
property.  Consequently every additive group homomorphism
\[
 h:X\longrightarrow\R
\]
from a separable real Banach space $X$ is continuous and hence $\R$-linear.
\end{lemma}

\begin{proof}
The first statement is the standard Baire-isomorphism theorem: modulo meagre
sets, every uncountable Polish space is Baire-isomorphic to Baire space
$\omega^\omega$, which is homeomorphic to the irrationals.  Countable Polish
spaces cause no problem.  Thus the hypothesis for subsets of the reals
transfers to every Polish space.  The usual construction of this isomorphism
is countable, so $\DC$ suffices.

For completeness, let $h:X\to\R$ be additive and put
\[
 A_k=h^{-1}([k,k+1)),\qquad k\in\mathbb Z.
\]
These Baire-property sets cover the Baire space $X$, so some $A_k$ is
nonmeagre.  Pettis' category lemma gives a neighbourhood $U$ of zero contained
in $A_k-A_k$.  Hence $|h(x)|\le1$ on $U$.  For any $m\ge1$, $m^{-1}U$ is a
neighbourhood of zero and additivity gives $|h(x)|\le1/m$ there.  Thus $h$ is
continuous at zero.  It is rationally homogeneous by additivity; approximation
of a real scalar by rationals and continuity make it $\R$-linear.
\end{proof}

The invocation of Pettis is harmless in $\ZF+\DC$: its proof is the standard
countable Baire-category argument.  The same observation supplies the Baire
category theorem and uniform boundedness for separable Banach spaces.

\section{Proof of the characterization}

\begin{proof}[Proof of Theorem~\ref{thm:main} inside $(\ast)$]
Fix $n\ge2$ and a functional $F$ as in $D_n$.  Radial functions identify
$\mathcal S_0^n$ with the positive cone $\C_+$, and radial addition becomes
pointwise addition.

Separate additivity extends $F$ uniquely, coordinate by coordinate, to a
separately additive map
\[
 T:\C^n\longrightarrow\R.                                      \tag{1}
\]
This is just the Grothendieck-group construction.  Explicitly, write
$f=f^+-f^-$ and expand in each coordinate with alternating signs.  The result
does not depend on a chosen decomposition: if $u-v=u'-v'$ in $\C_+$, then
$u+v'=u'+v$, and cancellation plus additivity gives the same value.  This
construction uses no choice principle.

Fix all but one argument in (1).  The remaining section $h:\C\to\R$ is an
additive homomorphism between Polish groups.  Lemma~\ref{lem:pettis} makes it
continuous and $\R$-linear.  This holds in every coordinate and for every
choice of the other arguments.  Restricting back to $\C_+$ yields two facts
that were absent from the question's hypotheses:
\begin{enumerate}
\item $F$ is positively homogeneous in every coordinate;
\item $F$ is separately bounded on star bodies contained in the Euclidean
unit ball.
\end{enumerate}
Indeed, after the other arguments are fixed, the corresponding continuous
linear functional on $\C$ is bounded on its unit ball.

We may now apply the separately bounded branch of Theorem 3.4 of
Jim\'enez--Villanueva.  That theorem states precisely that a separately
additive, positively homogeneous, separately bounded $F$ which vanishes when
two arguments intersect only at the origin has a representation
\[
 F(L_1,\ldots,L_n)
   =\int_{\Sph}\rho_{L_1}(u)\cdots\rho_{L_n}(u)\,d\mu(u)          \tag{2}
\]
for a finite signed regular Borel measure $\mu$.  Rotation invariance of $F$
makes $\mu$ rotation invariant.  Uniqueness of normalized invariant measure
on the compact homogeneous sphere gives $\mu=a\sigma$ for some $a\in\R$.
Comparison with the normalization of $\Vt$ yields
\[
 F(L_1,\ldots,L_n)=na\,\Vt(L_1,\ldots,L_n),
\]
which is $D_n$.

There is no hidden use of full Choice in this last paragraph.  For the bounded
case, the source proof uses joint continuity of a finite multilinear form,
the Riesz representation theorem on the compact metric sphere, countable
regularity, and the continuous orthosymmetric $C(K)$ representation lemma.
These are all available under $\DC$: joint continuity follows from the
Baire-category/uniform-boundedness argument, and the measure constructions
are countable on a compact metrizable space.  Equivalently, (2) follows by
iterating the standard identity
$B(f,g)=B(fg,1)$ for continuous orthosymmetric bilinear forms on $C(K)$ and
then applying compact-metric Riesz representation.
\end{proof}

\section{Consistency strength and scope}

Solovay's classical model would already prove compatibility, assuming an
inaccessible cardinal.  The Baire-property input is substantially cheaper.
Shelah proves in Conclusion 7.17 of \emph{Can you take Solovay's inaccessible
away?} that the following theories are equiconsistent:
\[
 \mathsf{ZFC}\qquad\text{and}\qquad
 \ZF+\DC+\text{``every set of reals has the Baire property''.}  \tag{3}
\]
Moreover $\ZF$ and $\mathsf{ZFC}$ are equiconsistent, by passing to the
constructible inner model.  Thus (3) and the internal proof establish
Theorem~\ref{thm:main} from $\operatorname{Con}(\ZF)$.

\begin{remark}[Exact logical scope]
This proves a full affirmative answer to the compatibility question.  It does
not prove $D_n$ in bare $\ZF$, nor does it prove independence.  The Axiom of
Choice gives the known counterexample $F=h\circ\Vt$ from a nonlinear Cauchy
additive $h:\R\to\R$, whereas the Baire-property model eliminates not only
such $h$ but every discontinuous additive section on $\C$.
\end{remark}

\begin{remark}[Why scalar Cauchy regularity alone is not the proof]
Knowing only that additive maps $\R\to\R$ are linear controls scalar rays but
does not itself bound a functional on the infinite-dimensional space $\C$.
The all-sets-Baire-property hypothesis is what makes every section
$\C\to\R$ continuous, exactly supplying the source theorem's missing
hypothesis.
\end{remark}

\section*{Novelty check}

Bounded exact-title, exact-question, phrase, author, citation, and Solovay-model
searches on 11 August 2026 found the two source papers and later work on other
dual-volume questions, but no later resolution of Question 3.5.  The argument
is short enough that it may be folklore; novelty confidence is therefore
moderate pending expert review.

\begin{thebibliography}{9}
\bibitem{JV}
C. H. Jim\'enez and I. Villanueva,
\emph{Characterization of dual mixed volumes via polymeasures},
J. Math. Anal. Appl. 426 (2015), 688--699;
\href{https://arxiv.org/abs/1408.6796}{arXiv:1408.6796}.

\bibitem{DGP}
P. Dulio, R. J. Gardner, and C. Peri,
\emph{Characterizing the dual mixed volume via additive functionals},
Indiana Univ. Math. J. 65 (2016), 69--91;
\href{https://arxiv.org/abs/1312.4072}{arXiv:1312.4072}.

\bibitem{Shelah}
S. Shelah,
\emph{Can you take Solovay's inaccessible away?},
Israel J. Math. 48 (1984), 1--47;
\href{https://doi.org/10.1007/BF02760522}{doi:10.1007/BF02760522}.

\bibitem{Solovay}
R. M. Solovay,
\emph{A model of set-theory in which every set of reals is Lebesgue
measurable}, Ann. of Math. 92 (1970), 1--56.
\end{thebibliography}

\end{document}
